\section{The evaluation that would actually decide the question}
\label{sec:evaluation}

Nothing in Section~\ref{sec:experiments} measures whether \forge{} would be a
better Lean tactic. This section specifies the experiment that would, in enough
detail that it could be pre-registered.

\subsection{Freeze the comparison, including its theorem knowledge}

Use a pinned Lean release, a compatible pinned Mathlib revision, fixed imported
modules, and recorded backend versions. Preserve each source goal with its
original hypotheses. Run methods in isolated processes so that one theorem's
auxiliary declarations and caches cannot leak into another method's state, and
reset proof state between runs so that a successful trial cannot add a lemma
that helps a later baseline run.

\emph{A fair test must specify both the theorem statement and the environment in
which it is presented.} This is the single most common way an automation
benchmark misleads, and it has a specific form here: a benchmark extracted from
an already compiled library can make the tested theorem itself, a renamed
duplicate, or a later theorem depending on it available through retrieval.
Renaming a theorem is not enough if its exact type remains indexed as a premise,
and attributes that indirectly reintroduce a held-out theorem are also leakage.
Build the retrieval index only from the \emph{pre-theorem} environment, or
generate each benchmark at the correct declaration point.

Never inspect the target theorem's existing proof to choose premises unless the
experiment is explicitly labelled an \emph{oracle-premise upper bound}.

\subsection{Baselines}

At minimum: stock \grind{}; tuned \grind{} with explicit budgets and relevant
suggestion settings; a simple public-tactic portfolio; and the full \forge{}
system. On applicable subsets add Aesop, \code{nlinarith}, \code{omega}, Duper,
the existing SOS and LP backends, \code{lean-smt}, and the shared-state
bit-vector route.

Two baselines matter more than the rest.

The first is the \emph{equal-knowledge} baseline. Use the same imports, local
context, theorem statement, typeclass environment, and premise universe. If
\forge{} performs premise retrieval, give the baseline a matched retrieval
condition or report the retrieval effect separately. Giving \forge{} an
associativity lemma while withholding the same available theorem from a
premise-aware baseline measures a data-access difference, not a search
improvement. A system that wins only against an artificially weak baseline has
demonstrated nothing.

The second is the \emph{non-communicating portfolio}: the same workers, the same
total budget, no fact sharing. This is the experiment that distinguishes
``integration earns its complexity'' from ``we added more solvers'', and it is
the one most easily omitted.

\subsection{Two tracks, three input conditions}

Run two tracks. The \emph{cold-knowledge} track allows only the actual
pre-theorem environment and tests missing-lemma synthesis. The
\emph{equal-premise} track gives all methods the same supplied premises and
tests proof search and cooperation.

For structural problems, distinguish three input conditions: definitions only;
a common relevant helper-lemma set; and an oracle-assisted set containing the
handwritten proof's helper lemmas. The gap between the first two measures lemma
acquisition; the third separates discovery failures from local proof failures.
A fourth ablation --- supply a hand-written correct induction motive to every
eligible method --- isolates motive selection from leaf solving.

\subsection{Corpus design}

Mix existing development goals with newly written examples, and stratify:

\begin{itemize}[leftmargin=1.4em]
\item recursive data-structure theorems needing induction or changing
  accumulators --- and not only append and reverse: include trees, mutually
  recursive syntax, nested recursion, and dependent vectors;
\item polynomial cost and invariant proofs;
\item semialgebraic inequalities with side conditions --- sparse and dense,
  near-zero and exactly-zero boundary behaviour, rationals with denominator side
  conditions, integer versus real distinctions;
\item quantified statements needing new arithmetic terms;
\item constructive existential specifications with scope-sensitive witnesses;
\item mixed guarded arithmetic and rewriting;
\item Boolean-heavy verification conditions --- theory-heavy and theory-light,
  not only pigeonholes;
\item higher-order and indexed-datatype obligations;
\item a large control set of goals \grind{} already handles well, including its
  arithmetic and bit-vector strengths.
\end{itemize}

Include valid but unsupported problems, false conjectures, misleading induction
variables, irrelevant assumptions, and goals requiring lemmas outside the first
retrieval batch. The control set is not optional: it detects whether the new
controller wastes resources on familiar cases.

Split by theorem dependency family or development module, not by random
near-duplicate statements. Hold out whole families when tuning heuristic
weights --- renaming variables or changing constants in a training problem does
not create an independent test set for a lemma-discovery heuristic. And do not
choose only problems generated from a favoured certificate language; that is
precisely the bias Section~\ref{sec:experiments} had to flag repeatedly in the
prototype evidence.

\subsection{Metrics}

Report solved goals at several fixed wall-clock and heartbeat budgets, together
with median and tail latency, peak memory, external-process time, proof and
certificate size, reconstruction time, and kernel-checking time. Distinguish
cold import costs from warm elaboration costs, and count any proof cache or
premise index that amortises work explicitly rather than charging it only to the
competitor.

Record separately: failed translations, unsupported fragments, timeouts,
rejected certificates, and residual reconstruction goals. \emph{Count a proof
with one residual goal as unsolved.} A solver verdict with an unclosed proof
hole is not a solved benchmark, and a worker's Boolean result is never the
benchmark's success oracle. Timeouts and reconstruction failures belong in the
denominator.

Keep deterministic replay success as its own metric. Report paired wins and
losses by problem family, not just total solved counts, using paired confidence
intervals or a paired test over an appropriate sampling unit --- projects or
theorem clusters, since examples within a module are correlated. For times, use
budget curves or censored-runtime summaries rather than averaging only the
successful cases and discarding timeouts.

Compatibility mode asks a different question and needs a different plot: how
much \emph{additional} resource rescues failures without replacing the
successful baseline run? Plot rescue coverage against that additional budget,
and never present the preservation property of
Proposition~\ref{prop:preservation} as an equal-cost performance theorem.

One class of outcome is not a metric at all. \emph{Any accepted false theorem or
unapproved axiom is a release-blocking failure, not an accuracy regression.}

\subsection{Required ablations}

Remove one capability at a time: backward obligation production; induction;
induction without generalisation; helper-lemma discovery; constrained
instantiation; constant-only witnesses; hidden-quadratic certificates;
hypothesis-product cones; Bernstein subdivision; zero-aware univariate search;
lattice witnesses; the Boolean lane; cross-worker fact exchange; and
relevance-based scheduling. Evaluate generalisation separately from
helper-lemma discovery so that the contribution of each is visible --- \src{p4}'s
$2\times2$ result (Table~\ref{tab:p4-ablation}) is the reason this matters.

And run the stateless-restart versus incremental-adapter comparison explicitly,
since it is the only thing that can justify the invasive integration work of
Gate~5.

\subsection{Acceptance gates}

A defensible acceptance criterion has three parts, and the threshold must be
chosen \emph{before} the run:

\begin{enumerate}[leftmargin=1.4em]
\item Retain a high fraction of the goals the pinned baseline solves.
\item Solve a substantial share of held-out families requiring combinations of
  induction, construction, nonlinear inequalities, and Boolean search.
\item Keep every result independently checkable, with proof reconstruction and
  kernel checking included in the cost.
\end{enumerate}

One contributing proposal offered, purely as an illustration of the \emph{form}
such a target takes, ``20\% more fully checked solves on the selected difficult
families and at most 5\% regression on the baseline suite''. That is a shape,
not a prediction: no such evaluation has been run, no percentage is forecast
here, and a 10-point gain in one category would not excuse an undisclosed
regression in the control category.

\subsection{A soundness regression suite, run continuously}

Separately from the performance benchmark, maintain a suite that tries to break
the system rather than to score it: adversarial certificates (oversized
exponents, duplicate keys, forward proof references, transplanted certificates,
omitted subdivision children, negative or float weights, truncated proof logs);
branch-leakage tests; circular-lemma tests; wrong-context replay; axiom-policy
violations; and compatibility checks against each supported Lean revision.

Mutate a denominator condition; swap two polynomial variables; replace a strict
inequality with a weak one; change an exponent; reverse an equality transport;
remove a branch from a box cover; alter a typeclass instance; point a proof node
to a future node; and replay under a different local context. Each of these has
an executable analogue in at least one prototype today, and each should have one
in the Lean implementation.

Resource exhaustion is a first-class outcome in this suite. A timeout or
out-of-memory event must return a recoverable status --- never trigger a
fallback that trusts the search output.
